%%%%%%%%%%%%%%%%%%%%%%% file template.tex %%%%%%%%%%%%%%%%%%%%%%%%%
%
% This is a general template file for the LaTeX package SVJour3
% for Springer journals.          Springer Heidelberg 2010/09/16
%
% Copy it to a new file with a new name and use it as the basis
% for your article. Delete % signs as needed.
%
% This template includes a few options for different layouts and
% content for various journals. Please consult a previous issue of
% your journal as needed.
%
%%%%%%%%%%%%%%%%%%%%%%%%%%%%%%%%%%%%%%%%%%%%%%%%%%%%%%%%%%%%%%%%%%%
%
% First comes an example EPS file -- just ignore it and
% proceed on the \documentclass line
% your LaTeX will extract the file if required
%\begin{filecontents*}{example.eps}
%!PS-Adobe-3.0 EPSF-3.0
%%BoundingBox: 19 19 221 221
%%CreationDate: Mon Sep 29 1997
%%Creator: programmed by hand (JK)
%%EndComments
%gsave
%newpath
%  20 20 moveto
%  20 220 lineto
%  220 220 lineto
%  220 20 lineto
%closepath
%2 setlinewidth
%gsave
%  .4 setgray fill
%grestore
%stroke
%grestore
%\end{filecontents*}
%
\RequirePackage{fix-cm}
%
%\documentclass{svjour3}                     % onecolumn (standard format)
%\documentclass[smallcondensed]{svjour3}     % onecolumn (ditto)
%\documentclass[smallextended]{svjour3}       % onecolumn (second format)
\documentclass[twocolumn]{svjour3}          % twocolumn
%
\smartqed  % flush right qed marks, e.g. at end of proof
%
\usepackage{graphicx}
\let\proof\relax
\let\endproof\relax
\usepackage{amsthm}

\usepackage{xcolor}

\usepackage{multirow}
\usepackage{makecell}
\usepackage{amsmath}
\theoremstyle{plain}
\newtheorem{myDef}{Definition}
\newcommand{\jf}[1]{\textcolor{purple}{#1}}
\let\Bbbk\relax
\usepackage{amssymb}
\usepackage{mathtools}

\usepackage{url}
\usepackage{booktabs}
\usepackage{graphicx}
\usepackage{marvosym}

\usepackage{subcaption}
\usepackage{enumitem}

\usepackage{balance}

\usepackage{ulem}
\usepackage[ruled,vlined,linesnumbered,ruled]{algorithm2e}

\newcommand{\xv}{\boldsymbol{x}}
\newcommand{\Xv}{\boldsymbol{X}}
\newcommand{\sv}{\boldsymbol{s}}
\newcommand{\zv}{\boldsymbol{z}}
\newcommand{\Zv}{\boldsymbol{Z}}
\newcommand{\Zvi}{\boldsymbol{Z}_i}

\newcommand{\eat}[1]{}
%\newcommand{\zy}[1]{\textcolor{blue}{#1}}
\newcommand{\zy}[1]{#1}

%\newcommand{\rev}[1]{\textcolor{blue}{#1}}
\newcommand{\rev}[1]{#1}
\newcommand{\rc}[1]{\vspace{3ex}\noindent\textbf{\textit{#1}}\vspace{2ex}}

\newcommand{\major}[1]{\textcolor{blue}{#1}}

\usepackage{framed}

\usepackage{xr}
\makeatletter
\newcommand*{\addFileDependency}[1]{% argument=file name and extension
\typeout{(#1)}
\@addtofilelist{#1}
\IfFileExists{#1}{}{\typeout{No file #1.}}
}\makeatother

\newcommand*{\myexternaldocument}[1]{
\externaldocument{#1}
\addFileDependency{#1.tex}
\addFileDependency{#1.aux}
}
\myexternaldocument{main}

%
% \usepackage{mathptmx}      % use Times fonts if available on your TeX system
%
% insert here the call for the packages your document requires
%\usepackage{latexsym}
% etc.
%
% please place your own definitions here and don't use \def but
% \newcommand{}{}
%
% Insert the name of "your journal" with
% \journalname{myjournal}
%
\begin{document}
\begin{sloppypar}
\title{Response to Review Comments of ``Towards GPU Memory-Aware Efficient Contrastive Shapelet Learning for Unsupervised Representation Learning in Multivariate Time Series''%\thanks{Grants or other notes
%about the article that should go on the front page should be
%placed here. General acknowledgments should be placed at the end of the article.}
}
%\subtitle{Do you have a subtitle?\\ If so, write it here}

\titlerunning{Response to Review Comments}        % if too long for running head

\author{Zhiyu Liang$^1$         \and
        Xiongfei Li$^1$ \and Chen Liang$^1$ \and Yuchen Hu$^1$ \and Yutong Zhao$^1$ \and Zheng Liang$^1$ \and Jianfeng Zhang$^2$ \and Hongzhi Wang$^{1}$ \and Lujia Pan$^2$   %etc.
}

%\authorrunning{Short form of author list} % if too long for running head

\institute{Zhiyu Liang \at
              % first address \\
              % Tel.: +123-45-678910\\
              % Fax: +123-45-678910\\
              \email{zyliang@hit.edu.cn} 
\and Hongzhi Wang (\textit{Corresponding author}) \at \email{wangzh@hit.edu.cn}\and
$^1$ Harbin Institute of Technology, Harbin, China\\
$^2$ Huawei Noah's Ark Lab, Shenzhen, China
}

\date{Received: date / Accepted: date}
% The correct dates will be entered by the editor


\maketitle

\section{List of changes}

First of all, we express our deepest gratitude to the reviewers for their insightful
comments on our paper. These comments have been fully taken into
account in this revision, and every effort has been made to address the comments
raised by the reviewers. The major changes are marked in \major{blue} in the
revised manuscript. In particular, following the suggestion of the editor, we omit the evaluation of large scale GPU clusters as pointed by O4 and explain the reasons in detail to the reviewer. We not only discuss the three methods recommended by O3 in our related work, but also provide additional experiments to compare them with our CSL using 30 classification tasks.

The main changes are summarized as follows.

\begin{enumerate}
%\vspace{-1.5ex}
    \item \textbf{Clarification of design choice and contributions.}  We reveal the principle behind why shapelet-based representation is better than general neural networks for time series (see \S~1.1). We further clarify the coherence of MAST with the original CSL (see \S~1.2) and the customized yet substantial nature of MAST (see \S~5).  

    \item \textbf{Additional evaluation and analysis.} We compare our CSL with three recently proposed baseline methods recommended by Reviewer \#1 (see \S~6.10), report and analyze the inference time of different URL methods (see \S~6.8) and the time of the profiling epoch (see \S~6.11.3). We also provide more discussion on the substantial improvement of our MAST planner (see \S~6.11.1), a systematic guide on how to configure CSL for better performance (see \S~6.4), the reason why CSL performs better with limited annotations (see \S~6.9), and the failure mode of CSL for high-dimensional data (see \S~6.2.1). Moreover, we reorganize the evaluation with limited annotations section as an extended evaluation, removing it from the main results to an additional subsection (see \S~6.9).

    \item \textbf{Clarification of implementation details and key terms.} We provide implementation details of the profiling and discuss the observed variation (see \S~5.2). We clarify that checkpointing specifically refers to activation checkpointing to avoid misunderstanding (see Abstract and \S~1.2). We also clarify that it is the BLP with the determined constants that can be solved precisely (see \S~5.2), and refine the terminology for describing a planner by replacing absolute claims of ``optimal" with more accurate descriptions like ``effective" and ``high-quality" (see \S~1.2, \S~2.3, and \S~5.2).

    \item \textbf{Update of related work and code repository.} We further discuss the methods recommended by Reviewer \#1 in the literature review (see \S~2.1) and update the repository to include the code of MAST (see the Github link in the abstract).
\end{enumerate}
\vspace{-0.5ex}

In the rest of this letter, we explain \textbf{point-by-point} how the concerns raised by the reviewers are addressed in the revised paper. The original Review Comments (\textbf{\textit{RC}}) are formatted in \textbf{\textit{italic and bold text}}, while our responses appear in Normal format and the manuscript text is shown in the \fbox{frame box}. We hope that these changes satisfy the reviewer's requirements, which will make the paper acceptable for publication.


\section{Response to Reviewer \#1}

\rc{O1. The methodological grounding is not convincing. Although the paper includes ablations and sensitivity studies, the design-space analysis remains limited, and the support for shapelet-based representation is primarily empirical rather than theoretical. In particular, why shapelet-based approach is better than others? Are there any other mechenisms that are even better than shapelet-based representation itself for this problem? 
The results show that shapelets work well within the proposed framework, but they do not convincingly establish why this representation should be preferred more broadly over other modern time-series SSL alternatives. In addition, the framework still appears fairly parameterized, with many design choices and no clearly systematic configuration methodology (how will a user configure those parameters for high performance? ).  As a result, the fundation of the paper is shaky.}

We thank the reviewer for the insightful feedback. Despite that we have clearly stated in the motivation part in \S 1.1 the limitations of alternative time-series SSL methods and why we intend to solve these problems by exploring shapelet-based representation, we are aware that our current analysis for choosing shapelet is more from empirical evidence rather than principles. In fact, the reason why shapelet-based representation is more effective is that it does not require learning the inherent structure of time series implicitly from a huge hypothesis space like using general neural networks. Instead, since shapelet-based features explicitly represent time series using the prototypical subsequences that have been widely shown superior for downstream time series tasks~\cite{Shapenet,liang2021efficient,liang2024fedst,USSL,MUSLA}, using it for URL can reduce the effective hypothesis space and thus simplify the optimization process to improve generalization. We provide more discussion in the revised paper as follows.

\begin{framed}
    \textbf{[Section 1.1]} 
    
    To address the issues mentioned above, our core idea is to explore the time-series-specific representation encoder without strong assumptions for URL. In particular, we consider the encoder based on a non-parametric time series analysis concept named \textit{shapelet}~\cite{ye2011time}, i.e. salient subsequence which is tailored to \rev{extract time series features from only important time windows to avoid the noises outside}. \major{Unlike general neural networks that require learning the inherent structure of time series implicitly from a huge hypothesis space, shapelet-based features reduce the effective hypothesis space by explicitly representing time series using prototypical subsequences, which has been widely shown superior for time series in tasks such as  classification~\cite{Shapenet,liang2021efficient,liang2024fedst} and clustering~\cite{USSL,MUSLA}. As a result, URL in such an effective and reduced hypothesis space can simplify the optimization process to improve generalization.} In addition, compared to the feature extracted from general neural networks, the shapelet-based feature can be more intuitive to understand for time series~\cite{ye2011time,liang2024fedst}.
\end{framed}


In addition, despite the fact that any deep URL method has a few hyper-parameters to tune and we have extensively shown that our CSL model is not very sensitive to parameter settings, we further provide a guidance for CSL configuration to better assist users in deploying our method in their new applications. The revision is as follows.

\begin{framed}
\textbf{[Section~\ref{subsec:exp:sens_analysis}]}

\noindent
\ref{subsec:exp:sens_analysis} Sensitivity analysis

We perform sensitivity analysis to study the key parameters, including the number of shapelet scales $R$ (default 8), the  minimum and maximum lengths of the shapelets $L_{min}$ (default $0.1T$) and $L_{max}$ (default $0.8T$), the decay rate $\alpha$ (default 0.5), and the regularization coefficients $\lambda$ (default 0.01) and $\lambda_S$ (default 1). 

......

\major{\textit{Guidance for CSL configuration.} Based on the above analysis and the discussion in \S~\ref{sec:CSL}, we provide a systematic guide for configuring our CSL in real-world applications. We recommend $R\!=8$, $L_{min}\!=\!0.1T$, $L_{max}\!=\!0.8T$, $\alpha\!=\!0.5$, $\lambda\!=\!0.01$, and $\lambda_S\!=\!1$ as our default for fast deployment and generally competitive performance across different scenarios. In cases users have expertise on pattern scales of their specific applications, they can better model these patterns by easily tuning $L_{min}$ and $L_{max}$. One can also determine the settings more precisely using hyper-parameter optimization tools, trading more tuning effort with higher representation quality.}
\end{framed}


\rc{O2. My other major concern is the paper's overall framing. The submission combines a shapelet-based URL method and a model-specific checkpointing planner, but I was not fully convinced that these two parts come together as one coherent journal-level contribution. In particular, MAST reads more like a supporting system optimization for CSL than a fully co-equal scientific contribution, which leaves the paper somewhat split in focus.}

We thank the reviewer for this insightful comment. We agree that the current presentation may give the impression that MAST is somewhat split from CSL. Our intention in the journal paper, however, is to address a \textbf{unified} problem: \textit{enabling efficient shapelet-based unsupervised representation learning under practical GPU memory constraints}. Specifically, \texttt{CSL} defines the representation mechanism, model architecture, and learning objective, while \texttt{MAST} is not merely an implementation detail but an integral component that bridges CSL algorithmic design and system efficiency. MAST enables CSL achieve state-of-the-art training efficiency under limited GPU memory footprint or state-of-the-art GPU memory optimization with the same time overhead, improving applications of the superior URL algorithm in realistic settings. We further clarify the coherence of the two contributions in our revised paper as follows.

\begin{framed}
    \textbf{[Section 1.2]}
    
To address C4, we propose \textsf{MAST}, the first customized checkpointing planner for our shapelet-based encoder. Unlike existing general planners, which decide to drop and recompute activations for all computational layers (e.g., convolution layer and BatchNorm layer) that usually have complex data dependencies~\cite{jain2020checkmate}, MAST determines the checkpointing plan in a more precise space by storing or dropping the shapelet-based sub-modules of different scales and measures, since they dominate the GPU memory consumption and can be recomputed independently. To search for the \major{high-quality} plan that maximally reduces training cost under a given memory budget, we theoretically analyze the computational cost and memory consumption with respect to the checkpointing plan. Then, we formulate the planning problem as a binary linear program that can be solved precisely with standard solvers. \major{It should be noted that MAST serves as an integral component that bridges CSL algorithmic design and system efficiency, enabling efficient shapelet-based URL under limited GPU memory, which is essential to enhance the practical applications of the superior CSL approach.}    
\end{framed}


\rc{O3. The evaluation is problematic and missing state-of-the-art baselins. Some baseline results are reproduced by the authors while others are taken from prior work, which makes the comparison less clean. The authors should reproduce all the results in the same testbad for fare comparision. The baseline suite also misses several recent time-series SSL methods, such as SimMTM (NeurIPS 2023), Ti-MAE (ICLR 2023), and SoftCLT (ICLR 2024), and the anomaly-detection study is less convincing than the classification and clustering parts due to its smaller scale and author-adapted baselines.}

We thank the reviewer for the constructive feedback. First, as also mentioned by the editor, it is not easy to reproduce all compared methods, especially for our experiments that include a total of 18 customized baseline methods originally implemented in different programming languages and some are even not open-sourced. However, the results taken from existing papers are only the classification accuracy of the baselines and the clustering RI and NMI of the tailored methods. On the one hand, such precision results are well-optimized by the original authors, which can be competitive enough for us to compare to evaluate our approach. On the other hand, there are still many precision results that come from our reproduction, while all efficiency results are run by ourselves on the same testbed to fully eliminate the influence of software and hardware environment. 

Second, we argue that the impression that anomaly detection study is less convincing than the classification and clustering parts is just because we conduct \textbf{much more experiments than standard} URL or SSL evaluation in the classification and clustering parts. In fact, since the target problem is a general-purpose URL for MTS, the standard and main evaluation is to compare CSL with existing URL approaches in different tasks and settings~\cite{TS-TCC,softclt}. However, in addition to the standard evaluation, we extensively compare many baseline methods specifically designed for specific tasks, which is optional but we believe useful. Moreover, as we have explained in \S~\ref{sec:setup}, most tailored anomaly detection methods predict anomalies for each data point rather than segment, and thus it is not very direct to fit them to the segment-level anomaly detection settings for a fair comparison, while the isolation forest used in this paper is a strong baseline in time-series anomaly detection.  

Finally, we further discuss the three recommended methods (where SimMTM has been discussed in the original manuscript) and provide more comparison with them in our experiments, which are released during the same period or later than our original CSL method. The overall results across all 30 classification tasks on the UEA datasets demonstrate that CSL outperforms all compared methods in terms of average accuracy, average ranking, and number of wins. The $p$-value in the accuracy difference from ours is 0.0727 for SoftCLT and less than 0.0034 for the others. This indicates that CSL is still competitive with the design based on the time-series-native shapelet. We also mention that Ti-MAE does not appear to be an officially published work in ICLR 2023. Instead, we find that it was submitted but rejected by ICLR 2023 and has not yet been published. Thus, we cite their preprinted version. 

The discussion and evaluation presented in our revision are as follows.

\renewcommand{\thetable}{\ref{tab:more_baselines}}
\begin{table}[t]
    \centering
    \caption{\major{Comparison with more recent URL baselines.}}
    \vspace{-2ex}
    \resizebox{\linewidth}{!}{\begin{tabular}{lcccc}
    \toprule
      \textbf{Statistic}   & SimMTM & Ti-MAE & SoftCLT & \textbf{CSL (Ours)} \\
    \midrule
       Avg Acc/AR  & 0.654/3.33	& 0.681/2.82 &	0.707/2.15 & \textbf{0.735}/\textbf{1.70}\\
    %\hline
    %   Avg Ranking & 3.33 & 2.93 & 2.05  &  \textbf{1.68}  \\
    \hline
       \makecell[l]{\eat{\text{1 vs. 1 Compasion} \\ (Wins/Ties/Loses)}W/T/L} & 25/2/3 & 23/1/6 & 19/1/10 & /\\
     \hline
    \text{$p$-val} &  \underline{0.0001} & \underline{0.0034} & 0.0727 & /\\
    \bottomrule
    \end{tabular}}
 %   \label{tab:more_baselines}
    \vspace{-3ex}
\end{table}

\begin{framed}
\textbf{[Section~\ref{subsec:exp:more_baselines}]}

\noindent
\ref{subsec:exp:more_baselines} \major{Comparison with more recent baselines}

\major{We also compare CSL with representative URL methods released during the same period or later than ours~\cite{liang2023-csl-vldb}, including SimMTM~\cite{dong2023simmtm}, Ti-MAE~\cite{li2023ti}, and SoftCLT~\cite{softclt}. The overall results across all 30 classification tasks on the UEA datasets are shown in Table~\ref{tab:more_baselines}. CSL outperforms all compared methods in terms of average accuracy, average ranking, and number of wins. The $p$-value in the accuracy difference from ours is 0.0727 for SoftCLT and less than 0.0034 for the others. This indicates that CSL is still competitive with the design based on the time-series-native shapelet.}      


    \textbf{[Section~\ref{subsec:survey_URL}]}

\noindent
    \ref{subsec:survey_URL} Unsupervised MTS representation learning 
    
Unlike in domains such as vision~\cite{SimCLR,PCL,yang2022vision} and natural language~\cite{SimCSE,nlp_repr_zhang2022unsupervised}, the study of URL in time series is still in its infancy. Generally, existing solutions can be divided into three categories, i.e., \textit{contrastive learning-based}, \textit{generative learning-based}, and \textit{hybrid} methods.  


Contrastive learning aims to learn meaningful representations by encouraging the encoder to distinguish between similar (positive) and dissimilar (negative) pairs of data samples, which has been widely shown to be effective in URL~\cite{SimCLR,SimCSE}. Inspired by word representation~\cite{word2vec}, T-Loss~\cite{T-Loss} adapts the triplet loss originally designed for natural language to achieve time series URL by contrasting subsequences that are close and distant in time steps. Similarly, BTSF~\cite{BTSF} uses the triplet loss to contrast different time series samples and combines the temporal and spectral features of MTS to enrich the representation. TNC~\cite{TNC} assumes semantic consistency between overlapping segments and contrasts such temporal neighborhoods and distant segments to model dynamic temporal patterns.  Instead of data representation in a single level, TS2Vec~\cite{ts2vec} combines timestamp-level with contextual series-level contrast for hierarchical time series representation, \major{while SoftCLT~\cite{softclt} makes a further improvement with soft assignment on the hierarchical contrastive loss.}  

Generative learning is an alternative way to achieve URL. Such methods try to learn the inherent structure of MTS by encouraging the representations to reconstruct or predict the time series (featuers) maximally. Autoencoding~\cite{baldi2012autoencoders} is a basic and popular generative method, which trains an encoder and a decoder simultaneously to embed raw data into representations and decode the representations to reconstruct the raw data, respectively. In recent years, masked pre-training in conjunction with Transformer~\cite{attention}, which has achieved great success in language modeling~\cite{kenton2019bert}, has become a potential generative paradigm for URL. The basic idea is to mask a portion of the input data and learn the representations that can correctly predict the masked parts using the unmasked data. TimeMAE~\cite{cheng2023timemae} \major{and Ti-MAE~\cite{li2023ti}} extend the masked autoencoder (MAE)~\cite{he2022masked-autoencoder} originally designed for image representation to time series. TST~\cite{TST} develops a general Transformer-based URL framework with a carefully designed masking strategy for MTS, showing advanced performance in several downstream analysis tasks.

Hybrid URL solutions combine contrastive and generative techniques to benefit from both. TimeDRL~\cite{chang2024timedrl} learns the representations of MTS by contrasting positive and negative samples and reconstructing the original data simultaneously. \major{SimMTM~\cite{dong2023simmtm}} proposes reconstruction from multiple neighbor time series while contrasting masked and unmasked input. CPC~\cite{CPC} uses historical data to predict the future in a latent representation space and encourages correct prediction by minimizing contrastive loss. TS-TCC~\cite{TS-TCC} expands this idea to perform both temporal prediction and contextual contrast to further improve the representation quality.

Although existing methods have achieved improvements in representation quality, they still have limitations such as lack of intuition in encoder design and dependency on specific assumptions, as discussed in \S~\ref{subsec:motivation}. 
\end{framed}

\rc{O4. The system scalability is unclear. The scalability evaluation remains limited in scope. It is largely confined to one V100 GPU and does not show broader system-level scaling across large GPU clusters.}

We thank the reviewer for the constructive feedback. However, as we have clearly shown in both theoretical analysis and experiments, the CSL workload has fairly good scalability to considerably large-scale input, and the trends in increase of runtime with input size can be observed from existing evaluation. In particular, all experiments can be run in just one V100 GPU, even for long time series of length up to 236784, which indicates the nice properties of our CSL in terms of lightweight and efficiency. In addition to the original CSL, the newly proposed MAST planner in the journal version enables the CSL model to run efficiently under a given memory budget, further reducing reliance on large-scale GPUs. 

Although we understand that in practice, one can almost always use more or newer GPUs that allow a larger batch size to process very large-scale data more efficiently, such a case is usually a common and trivial engineering improvement. Even if there may exist non-trivial system-level optimization problem, it is out of the scope of this paper where we aim at designing a more effective URL algorithm and a checkpointing planner that allows the algorithm runs efficiently in a very practical scenario with limited memory budgets. 

Therefore, as the editor has also mentioned, it is not only difficult for us to access large scale GPU clusters, but also not necessary for the workload like CSL to consider such large scales.


\rc{O5. MAST is a useful addition, but its contribution is relatively narrow. It is specialized to the Shapelet Transformer structure, evaluated in a limited setting, and its gains over the strongest checkpointing baselines are positive but modest. I also think the paper should be more careful in its discussion of optimality, since exact optimization for the specialized binary linear program does not imply optimality for the broader checkpointing problem.}

We thank the reviewer for the insightful comment. We believe that the contribution and evaluation of MAST is substantial. First, as we have extensively evaluated, the CSL algorithm is superior in many different datasets, downstream tasks, and settings. Therefore, the problem of deploying CSL in constrained GPU memory is fundamental to enhance the practicality of the algorithm from a system point of view. By carefully analyzing the limitations of existing general checkpointing solutions, it is a strong motivation to tailor a planner for the CSL model. Although the proposed MAST is customized to the CSL model, it is still substantial by enhancing the practicality of Shapelet Transformer across its wide range of applications.

Second, we believe that we have made an extensive evaluation of the MAST planner. The experiments use four representative datasets and cover a wide ranges, including training cost under different memory budgets, real memory consumption under different budgets, real-time memory footprint during training, training time with respect to input sizes under limited memory, and overhead for planning. In our revision, we further report the runtime of the first training epoch, i.e., the profiling epoch, and analyze the profiling overhead incurred by planners.

Third, we would like to argue that the improvement of our MAST to the state-of-the-art planners is not marginal, because we are under very limited GPU memory budgets (8-20 GB) and can hardly make trade-off between time and space. Specifically, as we have clearly stated, the improvement in efficiency of 1.43\%-7.65\% is under the same memory budget and the memory saving of 6.4\%-21.9\% is under the same training time. This system-level optimization is achieved without adding any hardware resource or modifying the model architecture and algorithm process. In addition, since GPU time is expensive, the improvement is quite practical and meaningful, especially for iterative development or large-scale training data that require a long training time. 

\renewcommand{\thetable}{\ref{tab:solver_time}}
\begin{table}[t]
    \centering
    \caption{Planning time (in ms) averaged over budget.}
    \vspace{-1.5ex}
      \resizebox{\linewidth}{!}{\begin{tabular}{lcc|ccc}
    \toprule
    \multirow{2}{*}{\textbf{Dataset}}  & \multirow{2}{*}{\textbf{\major{1st epoch}}} &\textbf{Heuristics} & \multicolumn{3}{c}{\textbf{Optimization}}\\
    \cline{3-6}
  \rule{0pt}{3ex} & & Mimose   & Checkmate & Monet & \textbf{MAST (Ours)} \\
      \midrule
   Cr & \major{1.15s} & 0.55$^\dag\pm$0.27 & 242$\pm$71 & 142$\pm$13 & \textbf{56}$\pm$17\\
   DD & \major{6.57s} &  0.69$^\dag\pm$0.18 & 362$\pm$40& 124$\pm$15 & \textbf{66}$\pm$23\\
   MI &\major{138.34s} & 0.62$^\dag\pm$0.31 & 372$\pm$32& 184$\pm$26 &\textbf{84}$\pm$18 \\
   PE & \major{3.37s}  & 0.59$^\dag\pm$0.31 & 360$\pm$44& 150$\pm$29 & \textbf{88}$\pm$10 \\
   \bottomrule
    \end{tabular}}
  %  \label{tab:solver_time}
    \vspace{-4ex}
\end{table}

We carefully revise our manuscript to further clarify the above issues as follows.
\begin{framed}
    \textbf{[Section~\ref{sec:MAST}]}

    This section presents \textsf{MAST}, a \underline{m}emory-\underline{a}ware checkpointing planner for the \underline{S}hapelet \underline{T}ransformer model of CSL to achieve efficient training under limited GPU memory. \major{Recall that ST is a unified shapelet-based encoder general to all MTS. Thus, although customized for ST, MAST is still substantial by enhancing the practicality of ST across its wide range of applications.} We first introduce the preliminaries in \S~\ref{subsec:pre_ckp}, and then detail the MAST method in \S~\ref{subsec:mast_method}.

    \textbf{[Section 6.11.1]}
    Our MAST consistently outperforms the baseline methods in terms of training time, achieving 1.21$\times$-2.72$\times$ speedup over PyTorch. Compared to state-of-the-art planners Mimose, Checkmate and Monet, MAST reduces the time overhead by 1.43\%-7.65\% under the same budget, saving 6.4\%-21.9\% memory to train the CSL model with the same computational cost (measured by fitting the time-budget curves using quadratic polynomial). \major{Note that the gains are substantial, especially for large-scale training data and iterative development that require a long training time, because GPU time is expensive (e.g., \$2.48/hour for a 16 GB V100 GPU in Google Cloud\footnotemark[4].}   
    
     \textbf{[Section 6.11.3]}
     
     \major{Moreover, we study the runtime of the first epoch, which includes the total profiling overhead required by all planners. With all sub-modules recomputed, this epoch spends 1.2\%, 0.63\%, 0.58\% and 0.24\% of the total training time with the default 200 epochs for Cr, DD, MI, and PE that are given the minimum budget in Fig.~\ref{fig:real_budget}, respectively. Although the percentage of time spent on profiling and planning will increase with the total number of epochs decreases, using the planners are still faster than PyTorch since only a subset of model layers are recomputed from the second epoch on. Note that the profiling can also be pre-processed before training, which is needed only once and can be reused for the planning under different budgets. As such, only planning will incur additional overhead to the training, which is more efficient especially for a few epochs.}   
\end{framed}
\footnotetext[4]{https://cloud.google.com/compute/gpus-pricing}


In addition, we remove the claims regarding ``optimal'' and use words like ``effective'' and ``high-quality'' to describe the checkpointng plan. The detailed revision is as follows.

\begin{framed}
    \textbf{[Section 1.2]}

    Unfortunately, existing general planners, though applicable to arbitrary deep models, \major{are still limited in finding high-quality checkpointing plans} due to their heuristic strategy~\cite{sublinear,DTR,mimose} or complex non-linear programming formulation that can hardly be solved accurately~\cite{jain2020checkmate,monet}. To capitalize on the superior URL performance of CSL, it is essential yet still challenging to determine a proper plan for the CSL encoder to achieve efficient training within the memory budget (\textbf{C4}).   

To address C4, we propose \textsf{MAST}, the first customized checkpointing planner for our shapelet-based encoder. Unlike existing general planners, which decide to drop and recompute activations for all computational layers (e.g., convolution layer and BatchNorm layer) that usually have complex data dependencies~\cite{jain2020checkmate}, MAST determines the checkpointing plan in a more precise space by storing or dropping the shapelet-based sub-modules of different scales and measures, since they dominate the GPU memory consumption and can be recomputed independently. To search for the \major{high-quality} plan that maximally reduces training cost under a given memory budget, we theoretically analyze the computational cost and memory consumption with respect to the checkpointing plan......

\textbf{[Section 2.3]}

Since checkpointing trades training efficiency for memory by recomputing the dropped intermediate results, a challenging yet essential problem is how to make a proper checkpointing plan given a limited memory budget to reduce training overhead. Existing planners are designed mainly for general model training. To ensure compatibility with all possible computation dependencies between different model layers, some work~\cite{sublinear,mimose,DTR} designs heuristics to greedily choose model layers for checkpointing based on their memory usage and recomputation cost. \major{Such solutions are easy to generalize but may ignore global information of the planning task.} Alternative studies~\cite{jain2020checkmate,monet} solve the problem in a more general way, modeling it as an optimization problem to minimize training overhead with memory consumption constraints. However, the formulated programs contain non-linear constraints (e.g., polynomial~\cite{jain2020checkmate} or quadratic~\cite{monet}) that cannot be solved \major{precisely, limiting the quality of the solutions.}  

\textbf{[Section~\ref{subsec:mast_method}]}

Based on the above observations, the \textit{core idea} of MAST is to search for \major{an effective} checkpointing plan by \textit{determining whether the intermediate outputs of each sub-module which dominates the GPU memory and can be recomputed independently are dropped} during the forward pass, while operations for deriving the loss are executed without checkpointing. In this way, we can model the planning problem as a binary linear program that can be solved precisely by standard solvers......
\end{framed}

\rc{O6. While the paper studies training efficiency and memory behavior, it says little about inference or processing latency in practical deployment settings. Given the paper's emphasis on downstream applicability, I would have liked to see at least some discussion or evaluation of this aspect.}

We thank the reviewer for the constructive suggestion. We add experimental results and analysis of the inference latency of CSL and the baseline URL methods in the revised manuscript. Overall, we can conclude from the results that the proposed CSL achieves superior URL performance with competitive inference latency, further demonstrating its practicality in real-world development. The details are as follows.

\renewcommand{\thefigure}{\ref{fig:scalability_infer}}
\begin{figure}
    \centering
\subcaptionbox{\centering Time w.r.t. $N$.
    \label{subfig:time-N-infer}}
{\includegraphics[width=.32\linewidth]{figures/running_time/time_N_infer.pdf}}
\subcaptionbox{\centering Time w.r.t. $D$.
    \label{subfig:time-D-infer}}
{\includegraphics[width=.34\linewidth]{figures/running_time/time_D_infer.pdf}}
\subcaptionbox{\centering Time w.r.t. $T$.
    \label{subfig:time-T-infer}}
{\includegraphics[width=.315\linewidth]{figures/running_time/time_T_infer.pdf}}
\captionsetup[subfigure]{labelformat=empty}
\\ %\vspace{-0.5ex}
\subcaptionbox{
    \label{subfig:legend_time_infer}}\centering
{\includegraphics[width=\linewidth]{figures/running_time_aug/legend-aug.pdf}}
\vspace{-5.5ex}
    \caption{\major{Inference time versus input size ($N$), dimension ($D$) and series length ($T$) on InsectWingbeat, DuckDuckGeese and EigenWorms respectively.}}
   % \label{fig:scalability_infer}
    \vspace{-2.5ex}
\end{figure}

\begin{framed}
    \textbf{[Section~\ref{sec:exp:time}]}

    \major{We also evaluate the inference time with $N$, $D$, $T$ on the above three datasets, respectively. The time is almost unchanged for each method with $D$ increasing to a considerably large value (Fig.~\ref{subfig:time-D-infer}), since different dimensions can be processed highly in parallel by a GPU. Unlike training, TNC is efficient for a small $T$ (e.g. $T=30$ for IW in Fig.~\ref{subfig:time-N-infer} and $T=270$ for DD in Fig.~\ref{subfig:time-D-infer}), but slows dramatically with $T$ (Fig.~\ref{subfig:time-T-infer}), because it divides the series into short segments and sequentially encodes batches of them during inference. TST runs out-of-memory for large T. TS2Vec and T-Loss are generally slower while TS-TCC and CSL are fairly efficient, mainly due to the different overhead and hardware utilization of their encoder architectures. Overall, the proposed CSL achieves superior URL performance with competitive inference latency, further showing its practicality in real-world development.}
\end{framed}

\section{Response to Reviewer \#2}

\rc{M1. Terminology: "Checkpointing" must be clarified.
The term "checkpointing" is used throughout without qualification, creating a serious risk of confusion. In common usage checkpointing refers to saving model weights to disk for fault tolerance and portability — a completely different technique from what this paper implements. What this paper describes is specifically gradient checkpointing or activation checkpointing — a memory optimization technique that discards intermediate activations during the forward pass and recomputes them during backpropagation, trading computation for memory reduction. These two techniques are fundamentally different in purpose, mechanism, and application. Using the unqualified term throughout risks misleading readers unfamiliar with the specific deep learning systems literature this draws from.
}


We thank the reviewer for the insightful feedback. To avoid misunderstanding and for ease of explanation, we use ``activation checkpointing'' the first time the term occurs in the main text. At the same time, we clarify that ``checkpointing'' specifically
refers to “activation checkpointing" in this paper. Then, we use ``checkpointing'' in the rest of the paper.

As for the abstract, we use ``activation checkpointing'' the first time it occurs and abbreviate it as AC in the rest content.

The revision is as follows.
\begin{framed}
\textbf{[Abstract]}

\noindent
    \textbf{Abstract} Unsupervised representation learning (URL) has great potential for multivariate time series analytics due to its capability to learn generalizable representation for many downstream tasks without using expensive labels. However, existing URL methods usually adopt the models originally designed for other domains (e.g. vision and language) to encode the time series data and {rely on strong assumptions to design learning objectives, which limits their representation ability}. To address these problems, we propose a novel URL framework for multivariate time series that aims to learn shapelet-based representation specific to time series via contrastive learning. To the best of our knowledge, we are the first to explore shapelet-based embedding in general-purpose URL. To achieve this goal, we particularly design a unified shapelet-based encoder and a novel learning objective of multi-grained contrasting and multi-scale alignment, and develop a data augmentation library to improve the generalization. To further enable efficient training under constrained GPU memory based on the widely used \major{activation checkpointing (AC)} technique, considering the limitations of existing general \major{AC} planning methods, we propose the first memory-aware
\major{AC} planner customized for our shapelet-based model.  We conduct extensive experiments using tens of real-world datasets and evaluate our proposal on many downstream tasks, e.g., classification, clustering, anomaly detection, etc. The results demonstrate the superiority of our method against not only URL competitors, but also techniques specially designed for downstream tasks, and showcase the effectiveness of our \major{AC} planner in improving training efficiency with limited GPU memory. Our code has been made publicly available at https://github.com/real2fish/CSL.

    \textbf{[Section 1.2]}

    To enable training under a restricted GPU memory budget without modifying the model architecture or using additional hardware, an effective solution is \major{activation checkpointing}\footnotemark[2], which stores only a fraction of the intermediate outputs in memory and recomputes the rest when needed. Although such a solution can reduce the GPU memory footprint, it slows down training due to the additional overhead incurred by  recomputation~\cite{mimose}. Thus, to effectively trade computation for memory, a crucial problem is making a checkpointing plan to decide when and where to drop and recompute the intermediate results to improve training efficiency without running out of memory. 
\end{framed}

\footnotetext[2]{\major{In this paper, we also use ``checkpointing'' to specifically refer to ``activation checkpointing'' for ease of explanation.}}

\rc{M2. The profiling mechanism lacks necessary technical detail.
Section 5.2 states that hardware-dependent constants are obtained by "profiling the computational layers on the running GPU" during the first training epoch. However the manuscript provides no implementation detail about how this is actually done. The following questions are left unanswered:\\ \\
Which PyTorch APIs are used? Is timing based on torch.cuda.Event? Is memory measured via torch.cuda.memory\_allocated() or torch.cuda.max\_memory\_allocated()?\\
How is individual sub-module memory isolated from total GPU memory given that PyTorch's caching allocator does not immediately release memory when tensors are deleted?\\
How is workspace memory distinguished from activation memory during sub-module execution?\\
How stable are the profiled constants across different mini-batches and what variance is observed?\\
How does profiling account for CUDA's asynchronous execution model?\\ \\
These are not minor details — they determine whether the BLP is solving an accurately specified problem or an approximation. GPU memory measurement in PyTorch is inherently noisy due to caching allocator behavior, memory fragmentation, and context-dependent effects. If the profiled constants are approximate the optimality guarantee of the BLP solution is also approximate. The authors must acknowledge this explicitly, provide a complete technical description of the profiling implementation, and discuss what accuracy guarantees can be made about the profiled constants and the resulting memory budget constraint.
}

We thank the reviewer for the constructive suggestion. In our implementation, we measure the total cost of each profiled layer using \textit{time.perf\_counter()} and call \textit{torch.cuda.synchronize()} for explicit synchronization. Before a layer runs, we adopt \textit{torch.cuda.reset\_peak\_memory\_stats()} to clear previous peaks, use \textit{torch.cuda.synchronize()} for synchronization, and record the baseline memory using \textit{torch.cuda.memory\_allocated()}. Then we use \textit{torch.cuda.max\_memory\_allocated()} to record the peak during layer execution, and the workspace memory is the difference between the peak and the baseline.

It should be noted that in our formulation in Eqs. (22)-(27), the workspace memory includes the total memory footprint for computing the current layer and storing all its intermediate outputs, while the stored activations of the previous layers that have finished running are hardware-independent, which can be derived from the size and type of the input data. Therefore, in the formulated problem of MAST, we do not need to specifically distinguish the activations of the sub-modules from the profiled workspace memory.     


In our profiling, we observe a variance to mean ratio less than 0.23\% for the cost and a negligible variance (less than $5e^{-6}$ of the mean) for the workspace memory, which empirically ensures that our planner can satisfy the memory constraints while considerably improving training efficiency, as shown in our evaluation in \S~\ref{subsec:exp:MAST}.

We clarify the above issues in our revised paper, which is as follows.

\begin{framed}
    \textbf{[Section~\ref{subsec:mast_method}]}

To limit memory usage, we first analyze the peak memory consumption during each iteration, which changes over layers and steps (i.e., forward, backward, and recomputation). Let $p_{_{r,m}}$ and $p_{_{\mathcal{L}}}$ be the peak memory for computing the sub-module $f_{r,m}$ and the loss layer during the forward pass. The peak memory when one layer is executed consists of
three parts, i.e., the memory used to store the intermediate outputs of the previously computed layers, the \major{workspace memory for computing the current layer and storing all its intermediate results}, and the memory consumed by storing the model parameters.

.....

To obtain these hardware-dependent constants \major{\textit{on the fly with minimal memory consumption}}, we checkpoint all sub-modules and profile them during the first training epoch. \major{Specifically, we measure the total cost of each profiled layer using \textit{time.perf\_counter()} and call \textit{torch.cuda.synchronize()} for explicit synchronization. Before a layer runs, we adopt \textit{torch.cuda.reset\_peak\_memory\_stats()} to clear previous peaks, use \textit{torch.cuda.synchronize()} for synchronization, and record the baseline memory using \textit{torch.cuda.memory\_allocated()}. Then we use \textit{torch.cuda.max\_memory\_allocated()} to record the peak during layer execution, and the workspace memory is the difference between the peak and the baseline.} 

The profiled \textit{cost} and \textit{workspace memory} averaged over different mini batches are assigned to the corresponding constants. Finally, we solve the formulated linear program and apply the identified plan to the remaining epochs. \major{We observe a variance to mean ratio less than 0.23\% for the cost and a negligible variance (less than $5e^{-6}$ of the mean) for the workspace memory, which empirically ensures that our planner can satisfy the memory constraints while considerably improving training efficiency, as shown in our evaluation in \S~\ref{subsec:exp:MAST}.}
\end{framed}

Second, our current discussion may give the impression that the BLP problem itself is accurate. In fact, since some constants of the BLP problem can be obtained only by profiling, which unavoidably introduces errors, the problem can never be absolutely accurate. Thus, our primary goal is to express that once the constants are determined, the BLP formulation can be solved precisely. We further clarify this issue in our revision as follows.

\begin{framed}
    \textbf{[Section~\ref{subsec:mast_method}]}

    Recall that the goal of our MAST is to minimize the training cost under the given memory budget. Based on the above analysis, we formulate the problem as \textit{optimizing the objective in Eq.~\eqref{eq:mast_objective} under the constraints of Eq.~\eqref{eq:mast_constraints}}. Note that all symbols except $\boldsymbol{y}$ in the optimization problem are constant. Thus, \major{the formulation is a binary linear program that can be precisely solved using standard solvers (e.g., GUROBI~\cite{gurobi}) once the constants are determined}. The hardware-independent constants $o_{r,m}$, $o_{_{
\mathcal{L}
}}$, $\theta$, $\Delta_{\mathcal{L}}$ and $\Delta_{r,m}$ can be easily derived using the size and type of the input data and the parameters (i.e. shapelets), while the constants $c_{r,m}$, $w_{r,m}$, $\tilde w_{r,m}$, $w_{_{\mathcal{L}}}$, $\tilde w_{_\mathcal{L}}$ can be obtained by profiling the computational layers on the running hardware. 
\end{framed}

In addition, we are aware that whether the profiled constants are absolutely accurate or not, it is difficult to achieve an optimal solution for the checkpointing planning problem itself (the formulated BLP is strictly not the same as the checkpointing planning problem). For a more precise expression, we remove the claims regarding ``optimal'' and use words like ``effective'' and ``high-quality'' to describe the checkpointng plan. The detailed revision is as follows.

\begin{framed}
    \textbf{[Section 1.2]}

    Unfortunately, existing general planners, though applicable to arbitrary deep models, \major{are still limited in finding high-quality checkpointing plans} due to their heuristic strategy~\cite{sublinear,DTR,mimose} or complex non-linear programming formulation that can hardly be solved accurately~\cite{jain2020checkmate,monet}. To capitalize on the superior URL performance of CSL, it is essential yet still challenging to determine a proper plan for the CSL encoder to achieve efficient training within the memory budget (\textbf{C4}).   

To address C4, we propose \textsf{MAST}, the first customized checkpointing planner for our shapelet-based encoder. Unlike existing general planners, which decide to drop and recompute activations for all computational layers (e.g., convolution layer and BatchNorm layer) that usually have complex data dependencies~\cite{jain2020checkmate}, MAST determines the checkpointing plan in a more precise space by storing or dropping the shapelet-based sub-modules of different scales and measures, since they dominate the GPU memory consumption and can be recomputed independently. To search for the \major{high-quality} plan that maximally reduces training cost under a given memory budget, we theoretically analyze the computational cost and memory consumption with respect to the checkpointing plan......

\textbf{[Section 2.3]}

Since checkpointing trades training efficiency for memory by recomputing the dropped intermediate results, a challenging yet essential problem is how to make a proper checkpointing plan given a limited memory budget to reduce training overhead. Existing planners are designed mainly for general model training. To ensure compatibility with all possible computation dependencies between different model layers, some work~\cite{sublinear,mimose,DTR} designs heuristics to greedily choose model layers for checkpointing based on their memory usage and recomputation cost. \major{Such solutions are easy to generalize but may ignore global information of the planning task.} Alternative studies~\cite{jain2020checkmate,monet} solve the problem in a more general way, modeling it as an optimization problem to minimize training overhead with memory consumption constraints. However, the formulated programs contain non-linear constraints (e.g., polynomial~\cite{jain2020checkmate} or quadratic~\cite{monet}) that cannot be solved \major{precisely, limiting the quality of the solutions.}  

\textbf{[Section~\ref{subsec:mast_method}]}

Based on the above observations, the \textit{core idea} of MAST is to search for \major{an effective} checkpointing plan by \textit{determining whether the intermediate outputs of each sub-module which dominates the GPU memory and can be recomputed independently are dropped} during the forward pass, while operations for deriving the loss are executed without checkpointing. In this way, we can model the planning problem as a binary linear program that can be solved precisely by standard solvers......
\end{framed} 

\rc{M3. The overhead reporting is incomplete.
The paper claims planning overhead is "only 0.06\% of total training time." However this counts only the BLP solving time and ignores the profiling epoch itself, which constitutes one full additional training epoch. For a standard 200-epoch run the true overhead is approximately 0.5\%, not 0.06\%. While 0.5\% remains negligible this selective reporting is misleading. The authors should report complete overhead including the profiling epoch and discuss how this scales for shorter training runs.
}

We thank the reviewer for the insightful comment. In the original manuscript, we have clarified that the planning time specifically refers to the time to make the plan by executing the heuristics or solving the optimization problems of different planners. We believe that this evaluation is important to be assessed independently, because while the profiling overhead is shared by all planners and can be pre-processed before training starts, the planning overhead has to be incurred to training each time a new budget is given, and the overhead is different for different planning strategies. 

Based on the suggestion of the reviewer, we further report the time of the full profiling epoch, and discuss how the overhead scales with fewer training epochs. The revision is as follows.

\renewcommand{\thetable}{\ref{tab:solver_time}}
\begin{table}[t]
    \centering
    \caption{Planning time (in ms) averaged over budget.}
    \vspace{-1.5ex}
      \resizebox{\linewidth}{!}{\begin{tabular}{lcc|ccc}
    \toprule
    \multirow{2}{*}{\textbf{Dataset}}  & \multirow{2}{*}{\textbf{\major{1st epoch}}} &\textbf{Heuristics} & \multicolumn{3}{c}{\textbf{Optimization}}\\
    \cline{3-6}
  \rule{0pt}{3ex} & & Mimose   & Checkmate & Monet & \textbf{MAST (Ours)} \\
      \midrule
   Cr & \major{1.15s} & 0.55$^\dag\pm$0.27 & 242$\pm$71 & 142$\pm$13 & \textbf{56}$\pm$17\\
   DD & \major{6.57s} &  0.69$^\dag\pm$0.18 & 362$\pm$40& 124$\pm$15 & \textbf{66}$\pm$23\\
   MI &\major{138.34s} & 0.62$^\dag\pm$0.31 & 372$\pm$32& 184$\pm$26 &\textbf{84}$\pm$18 \\
   PE & \major{3.37s}  & 0.59$^\dag\pm$0.31 & 360$\pm$44& 150$\pm$29 & \textbf{88}$\pm$10 \\
   \bottomrule
    \end{tabular}}
   % \label{tab:solver_time}
    \vspace{-4ex}
\end{table}

\begin{framed}
    \textbf{[Section 6.11.3]}

\noindent
6.11.3 Study of planning overhead

Except for the training time by applying different checkpointing plans as analyzed above, we study the additional overhead incurred by the planners, i.e., the time spent in making the plan. As Table~\ref{tab:solver_time} shows, the planning time of Mimose (marked by $\dag$), which is based on simple greedy heuristics, is significantly less than that of the optimization-based solutions. By scheduling recomputation of different sub-modules rather than all computational layers, MAST is consistently much faster than the counterparts Checkmate and Monet, achieving a speedup of 1.70$\times$-5.48$\times$. In conclusion, the proposed MAST can make a better checkpointing plan with reasonable planning time. We also notice that the planning overhead is almost negligible compared to the training overhead because of the simple structure of our ST model. For example, even for the fastest-to-train dataset in Fig.~\ref{fig:real_budget} (i.e., Cr) with a small budget (i.e., 8 GB), the planning time of MAST is only 0.06\% of the total training time with the default 200 epochs. 

\major{Moreover, we study the runtime of the first epoch, which includes the total profiling overhead required by all planners. With all sub-modules recomputed, this epoch spends 1.2\%, 0.63\%, 0.58\% and 0.24\% of the total training time with the default 200 epochs for Cr, DD, MI, and PE that are given the minimum budget in Fig.~\ref{fig:real_budget}, respectively. Although the percentage of time spent on profiling and planning will increase with the total number of epochs decreases, using the planners are still faster than PyTorch since only a subset of model layers are recomputed from the second epoch on. Note that the profiling can also be pre-processed before training, which is needed only once and can be reused for the planning under different budgets. As such, only planning will incur additional overhead to the training, which is more efficient especially for a few epochs.}     
    
\end{framed}


\rc{M4. MAST's scope and applicability are overstated.
MAST applies exclusively to the CSL/Shapelet Transformer and cannot be transferred to other models without complete redesign. The claim in Section 1.2 that MAST "enables broad applicability across a wide range of scenarios" is misleading. MAST's clean BLP formulation is only possible because the Shapelet Transformer has an independence structure built into it by the authors' own prior design — this is not a general advance in gradient checkpointing methodology. All claims about broad generalizability must be removed or substantially qualified.
}

We thank the reviewer for the insightful suggestion. We would like to clarify that MAST is a planner customized for the CSL model (i.e., Shapelet Transformer), which we have repeatedly stressed in the manuscript. We are aware that the description in Section 5 (we do not find such claims in Section 1.2) that MAST "enables broad applicability across a wide range of scenarios" may give an impression that it is MAST that can be generalized to models other than Shapelet Transformer. However, our original discussion intends to express that it is CSL that is general to a wide range of scenarios, and thus the customized MAST can also fit in these scenarios by enhancing CSL. 

To further clarify our idea, we substantially qualify the expression in our revised paper as follows.

\begin{framed}
    \textbf{[Section 5]}

    This section presents \textsf{MAST}, a \underline{m}emory-\underline{a}ware checkpointing planner for the \underline{S}hapelet \underline{T}ransformer model of CSL to achieve efficient training under limited GPU memory. \major{Recall that ST is a unified shapelet-based encoder general to all MTS. Thus, although customized for ST, MAST is still substantial by enhancing the practicality of ST across its wide range of applications.} We first introduce the preliminaries in \S~\ref{subsec:pre_ckp}, and then detail the MAST method in \S~\ref{subsec:mast_method}.
\end{framed}


\rc{M5. The code repository does not include the MAST implementation.
The manuscript claims "Our code has been made publicly available at https://github.com/real2fish/CSL." However this repository reflects the original PVLDB 2023 paper only and has not been updated to include MAST.
}

We thank the reviewer for the kind suggestion. We have updated the repository to include the code of MAST. 

\rc{M6. The limited-label evaluation should be reframed as extended validation.
No new methodology is introduced — the existing CSL framework is applied to subsampled label sets. It is useful validation but should not be listed as a primary contribution. The authors should provide analytical insight into why CSL is robust to label sparsity rather than merely demonstrating that it is.
}

We thank the reviewer for the constructive feedback. We have separated the evaluation with limited annotations from the main result part and reframed it as an extended validation in Section 6.9 of the revised paper. In addition, we discuss the reason for CSL to achieve superior performance. The detailed revision is as follows.

\begin{framed}
\textbf{[Section 6.9]}

\noindent
6.9 \major{Evaluation with limited annotations}
To further quantify the representation quality of the URL methods, we evaluate the classifier trained using limited annotations. To this end, we randomly sample a subset of labeled training data with a ratio ranging from $\{1\%,5\%,10\%,20\%,50\%\}$ to train the SVM on top of the representation and report the test accuracy averaged over five resamples. To ensure both label sparsity and class balance, we consider the datasets that have more than 20 MTS for every class and sample data independently with the sample size rounded up for different classes. The results are shown in Table~\ref{tab:semi-supervise}, from which we have three key observations.

First, accuracy generally increases with the number of labeled MTS for all methods, showing a reasonable and practical tradeoff between classification performance and available annotations. Second, the URL methods, especially T-Loss, TS-TCC, and CSL, achieve comparable or better overall performance compared to applying SVM on raw data (denoted Raw), validating the effectiveness of URL in extracting informative MTS features that are more robust to label sparsity. Last but not least, CSL achieves the best overall performance, which has the highest average accuracy and ranking, outperforms all competitors in most datasets and ratios in one-versus-one comparison, and is statistically significant better than the compared methods in terms of accuracy with a $p$-value of no more than 0.002. \major{This can be explained because MTS of different classes in the CSL representation form more grouped and separated clusters (as illustrated in Fig.~\ref{fig:repr}), and thus the classification boundary can be identified easily with a limited samples of labels.} The results further indicate the superior URL performance of our proposal. 
\end{framed}

\rc{M7. The failure mode analysis is insufficient.
CSL's poor performance on DuckDuckGeese (D=1345) is noted in one sentence without investigation. It is needed to examine whether this is a fundamental limitation of shapelet-based encoding, a hyperparameter issue, or an architectural constraint.
}

We thank the reviewer for the insightful feedback. In fact, the reason for our claim that the poor performance of CSL in DuckDuckGeese may indicate its limitation for high-dimensional data is because, as formally described in Eq.~\eqref{eq:general_shapelet}, the (dis)similarity of all dimensions is simply summed in the design of general shapelet-based encoding. Such an encoding mechanism, although usually effective for different MTS as our extensive evaluation shows, can potentially overlook complex but important correlations between dimensions.

We provide more discussion on this failure mode in our revised paper. The details are as follows.

\begin{framed}
    \textbf{[Section 6.2.1]}

    ...... Furthermore, we observe that CSL performs poorly on DuckDuckGeese (DD), which has a very high dimension of 1345 (see Table~\ref{tab:uea_statistics}). \major{This may indicate a limitation of CSL in modeling high-dimensional MTS, since the simple sum of (dis)similarity of all dimensions as in Eq.~\eqref{eq:general_shapelet} can potentially overlook complex but important correlations between dimensions.}
\end{framed}

\section{Response to Reviewer \#3}

We thank the reviewer for the invaluable time and effort spent reviewing our manuscript and the high praise of our work. Following the kindly guide of the editor, we have made every effort to address all the concerns raised by the reviewers, which we believe will further improve the quality of the manuscript. 

%\normalem
%\bibliographystyle{spmpsci} 
%\bibliography{references}
\end{sloppypar}
\end{document}